\documentclass[11pt]{article}
\usepackage{acl}
\usepackage{times}
\usepackage{latexsym}
\usepackage[T1]{fontenc}
\usepackage[utf8]{inputenc}
\usepackage{microtype}
\usepackage{inconsolata}
\usepackage{graphicx}
\usepackage{booktabs}
\usepackage{amsmath}
\usepackage{amssymb}
\usepackage{multirow}
\usepackage{tipa}
\usepackage[section]{placeins}

\title{Low-Resource Character-Level Grapheme-to-Phoneme Conversion with an Encoder--Decoder Transformer}

\begin{document}
\maketitle

\begin{abstract}
This paper describes a character-level encoder--decoder Transformer for
grapheme-to-phoneme (G2P) conversion on English and Spanish. The system uses a
bidirectional encoder with grouped-query attention and rotary position
embeddings, and an autoregressive decoder with learned absolute positional
embeddings. Evaluation is conducted under strict train/test grapheme
deduplication, which prevents memorised pronunciations from inflating test
scores. Monolingual exact match is 0.153 for English and 0.620 for Spanish;
multilingual training improves these to 0.210 and 0.797 respectively. Error
analysis shows that English failures are concentrated in reduced vowels,
schwa-like alternations, and learned or loanword vocabulary. The results expose
a sharp memorisation/generalisation gap in low-resource English G2P, while
showing that multilingual training helps but does not erase the effect of
orthographic depth.
\end{abstract}

\section{Introduction}
\label{sec:introduction}

Grapheme-to-phoneme conversion maps written word forms to pronunciations and is
a central component in text-to-speech, automatic speech recognition, and lexical
resource construction \citep{cheng2024survey}. The task is simple to state but
typologically uneven: languages differ sharply in how transparently orthography
encodes pronunciation. English is difficult because spelling reflects multiple
historical layers, stress is usually not written, and vowel reduction is
conditioned by prosodic structure. Spanish provides a useful contrast because
its orthography is comparatively regular, although stress placement and
loanwords still create errors.

This contrast makes G2P a useful test case for orthographic depth. In a shallow
orthography, much of the target phone string can be inferred from local
letter--sound correspondences. In a deep orthography such as English, however,
the same spelling pattern can map to different vowels depending on stress,
morphology, lexical frequency, and etymological stratum. A low-resource
character-level model therefore faces two different problems at once: learning
productive spelling--sound correspondences, and deciding when the written form
is insufficient and lexical memorisation is required.

This work makes three contributions. First, it implements an encoder--decoder
Transformer for character-level G2P, extending a class decoder-only baseline
with a bidirectional source encoder. Second, it evaluates under strict
deduplication so that test graphemes are not present in training. Third, it
characterises the memorisation/generalisation gap in low-resource English G2P:
seen-word performance can be near-perfect, but unseen English words remain hard.

\section{Model}
\label{sec:model}

The model follows the encoder--decoder Transformer architecture
\citep{vaswani2017attention}. The source sequence is a sequence of grapheme
characters, optionally preceded by a language tag in multilingual training. The
target sequence is a whitespace-tokenised IPA phone sequence.

The encoder contains six bidirectional layers with model dimension
$d_{\text{model}}=384$, six query heads, and two key/value heads. Self-attention
therefore uses grouped-query attention (GQA), in which multiple query heads
share each key/value head \citep{ainslie2023gqa}. Rotary position embeddings
(RoPE) are applied to encoder queries and keys \citep{su2021roformer}. For a
two-dimensional subspace, RoPE rotates a vector by position-dependent angle
$m\theta_i$:
\begin{equation}
\label{eq:rope}
R_{\Theta,m}
\begin{bmatrix}x_{2i}\\x_{2i+1}\end{bmatrix}
=
\begin{bmatrix}
\cos(m\theta_i) & -\sin(m\theta_i)\\
\sin(m\theta_i) & \cos(m\theta_i)
\end{bmatrix}
\begin{bmatrix}x_{2i}\\x_{2i+1}\end{bmatrix}.
\end{equation}
Feed-forward layers use SwiGLU activations \citep{shazeer2020glu} with
$d_{\text{ff}}=1024$. Layers use pre-layer normalisation residual blocks and
stochastic depth increasing linearly from 0 to 0.1 across depth.

The decoder contains six autoregressive layers. Each layer applies causal
multi-head self-attention, cross-attention over the final encoder states, and a
SwiGLU feed-forward block. Decoder positions are learned absolute embeddings
rather than RoPE. This asymmetry is deliberate: the encoder benefits from
relative grapheme distance, while the decoder generates short phone sequences
where absolute output position is sufficient. All decoder layers attend to the
same final encoder output, with decoder queries and encoder keys/values in
cross-attention. The output projection is tied to the target embedding matrix,
following \citet{press2017using}. The resulting monolingual English and Spanish
models contain approximately 19M trainable parameters.

This positional asymmetry is also linguistically motivated. Grapheme strings
contain local and semi-local dependencies: digraphs, vowel contexts, suffixes,
and neighbouring consonants often determine phone choice, and these patterns may
occur at different absolute positions in words of different lengths. RoPE is
therefore useful in the encoder because it represents relative distance between
graphemes. The decoder, by contrast, produces a short, mostly monotonic phone
sequence in which output position is less structurally variable; learned
absolute positions are sufficient, while cross-attention supplies the
source-side information needed for alignment.
Table~\ref{tab:hyperparams} summarises the main architecture, training, and
decoding settings used in the reported experiments.

\begin{table}[!htbp]
\centering
\small
\begin{tabular}{ll}
\toprule
Setting & Value \\
\midrule
Model dimension & 384 \\
Encoder / decoder layers & 6 / 6 \\
Attention heads & 6 query heads, 2 KV heads \\
Feed-forward dimension & 1024 \\
Activation & SwiGLU \\
Layer normalisation & Pre-LN \\
Stochastic depth & 0--0.1 across depth \\
Label smoothing & 0.1 \\
Beam size / length penalty & 4 / 0.6 \\
Checkpoint averaging & top 3 validation checkpoints \\
\bottomrule
\end{tabular}
\caption{Main architecture, training, and decoding settings.}
\label{tab:hyperparams}
\end{table}

\section{Data and Experimental Setup}
\label{sec:setup}

Experiments use SIGMORPHON-style G2P data: English and Spanish each contain
2,400 training items, 300 validation items, and 300 test items. Inputs are
character-level grapheme sequences; outputs are whitespace-tokenised IPA phone
sequences.

The whitespace-tokenised IPA format mostly simplifies the task by making phones,
rather than raw Unicode characters, the prediction units. This avoids forcing
the model to learn that multi-character symbols such as affricates or
diacritically marked phones should behave as single phonological units. However,
Spanish exposed a practical tokenisation issue: stress markers were initially
split from the following vowel during decoding, while the reference format
expected them to be attached. This made exact match artificially zero despite
many otherwise correct pronunciations. The final decoder output is normalised
back to the reference convention by reattaching stress markers before scoring,
illustrating that G2P evaluation is highly sensitive to tokenisation
conventions.

A critical preprocessing issue is train/test overlap. Two strategies were
implemented. Strategy A, used for the main results, removes any training or
validation item whose grapheme appears in the test set. This is the honest
generalisation setting because a test word cannot be answered by memorised
lookup. Strategy B keeps overlap but tags test items as seen or unseen; it is
used only for memorisation analysis. In Strategy B, seen-word exact match was
0.958, while unseen-word exact match was 0.156, demonstrating that overlap can
mask poor generalisation.

Training uses AdamW with cosine warmup, peak learning rate $2\cdot10^{-4}$ or
$3\cdot10^{-4}$ depending on the run, warmup of 8 or 5 epochs, label smoothing
$\epsilon=0.1$, dropout, stochastic depth, and checkpoint averaging over the
top three validation checkpoints. The smoothed target distribution is:
\begin{equation}
\label{eq:ls}
q(k\mid y)=
\begin{cases}
1-\epsilon, & k=y,\\
\epsilon/(|V|-2), & k\neq y,\ k\neq \textsc{pad},\\
0, & k=\textsc{pad}.
\end{cases}
\end{equation}
Decoding uses beam search with beam size 4 and length penalty 0.6
\citep{wu2016gnmt}. Training-only grapheme augmentation applies random swap,
delete, and duplicate operations with probability 0.15, motivated by robust G2P
augmentation work \citep{zhao2022rg2p}.

Validation used two complementary signals. Validation loss was computed with
teacher forcing, so it measures how well the model predicts the next phone given
the gold previous phones. This is useful for checkpoint selection but can
underestimate exposure-bias errors because the model is not conditioned on its
own earlier mistakes. Validation PER and exact match were therefore computed
with greedy autoregressive decoding, while final test evaluation used beam
search.

\section{Results and Analysis}
\label{sec:results}

Phone error rate (PER) is computed as token-level Levenshtein distance
normalised by reference length:
\begin{equation}
\label{eq:per}
\mathrm{PER}=\frac{\sum_i d(\hat{\mathbf{y}}_i,\mathbf{y}_i)}
{\sum_i |\mathbf{y}_i|}.
\end{equation}

\begin{table}[!htbp]
\centering
\small
\begin{tabular}{llccc}
\toprule
System & Lang & EM & PER & WED \\
\midrule
Monolingual & EN & 0.153 & 0.389 & 2.221 \\
Monolingual & ES & 0.620 & 0.107 & 0.841 \\
Multilingual & EN & 0.210 & 0.298 & 2.087 \\
Multilingual & ES & 0.797 & 0.051 & 0.443 \\
\bottomrule
\end{tabular}
\caption{Strict-deduplication test results. Multilingual training improves both
languages, with the largest absolute gain for Spanish and a smaller but
meaningful reduction in English phone error rate.}
\label{tab:results}
\end{table}

Table~\ref{tab:results} shows a large English--Spanish gap. Spanish reaches
0.620 exact match and low PER under the same data scale and architecture. This
supports the role of orthographic regularity in G2P difficulty
\citep{cheng2024survey}. English, by contrast, achieves only 0.153 exact match
on unseen words. The model often predicts plausible-looking pronunciations, but
fails exact word-level transcription.

Multilingual training improves both languages, but asymmetrically. Spanish
rises from 0.620 to 0.797 exact match and halves PER from 0.107 to 0.051.
English also improves, with exact match rising from 0.153 to 0.210 and PER
falling from 0.389 to 0.298, but it remains much harder than Spanish. This
suggests that the shared model benefits from additional character-to-phone
evidence and cross-lingual regularities, while English stress-conditioned vowel
reduction and loanword strata remain underdetermined from spelling alone.

\begin{table}[!htbp]
\centering
\small
\begin{tabular}{lcc}
\toprule
Evaluation setting & Seen EM & Unseen EM \\
\midrule
Strategy B: overlap retained & 0.958 & 0.156 \\
Strategy A: strict deduplication & -- & 0.153 \\
\bottomrule
\end{tabular}
\caption{English memorisation and generalisation under the two evaluation
strategies. Strategy A is used for the main results because it removes
train/test grapheme overlap.}
\label{tab:memorisation}
\end{table}

The memorisation results in Table~\ref{tab:memorisation} explain why the
strict setting matters. When repeated graphemes are allowed, the model can
nearly solve seen-word G2P by storing lexical pronunciations. Once repeated
test graphemes are removed, English performance collapses to the unseen-word
level. This is not a failure of sequence modelling alone; it reflects the
amount of lexical and stress-conditioned information that a 2,400-item English
training set cannot supply.

\begin{figure}[!htbp]
\centering
\IfFileExists{../figures/training_curves.pdf}{%
\includegraphics[width=\linewidth]{../figures/training_curves.pdf}}{%
\fbox{\parbox{0.92\linewidth}{\centering Placeholder for English and Spanish
training curves generated from \texttt{history.json}.}}}
\caption{Validation loss and phone-error-rate curves for English and Spanish.
The historical Spanish exact-match log is not plotted because an earlier
stress-marker tokenisation error made that validation series unreliable; the
final strict test exact match is reported in Table~\ref{tab:results}.}
\label{fig:curves}
\end{figure}

Figure~\ref{fig:curves} shows Spanish stabilising at lower PER than English,
matching the final test gap and suggesting that English difficulty is not only
a decoding artefact.

English errors are linguistically structured. The most frequent substitutions
involve \textschwa{} vs.\ \textsci{}, \textschwa{} vs.\ \ae{}, and
\textepsilon{} vs.\ \textschwa{}, indicating
confusion among schwa, lax vowels, and stress-sensitive vowel qualities. This
is expected: English vowel reduction neutralises contrasts in unstressed
positions, and orthography generally does not encode the stress information
needed to predict it \citep{flemming2005vowel}. Uniform label smoothing may
exacerbate this by spreading probability mass over phonologically implausible
substitutions.
Table~\ref{tab:errors} gives representative high-frequency substitutions.

\begin{table}[!htbp]
\centering
\small
\begin{tabular}{lll}
\toprule
Reference & Prediction & Interpretation \\
\midrule
\textschwa{} & \textsci{} & reduced/lax vowel confusion \\
\textsci{} & \textschwa{} & unstressed vowel reduction \\
\textschwa{} & \ae{} & reduced vs full vowel \\
\textepsilon{} & \textschwa{} & stress-conditioned reduction \\
s & z & consonant voicing ambiguity \\
\bottomrule
\end{tabular}
\caption{Representative high-frequency English substitution errors. The
dominance of reduced-vowel errors shows that many failures are linguistic
rather than random decoding mistakes.}
\label{tab:errors}
\end{table}

The hardest English examples are also revealing: \textit{disopyramide},
\textit{oneiromancy}, \textit{impregnable}, \textit{jambeau}, and
\textit{remixture} are specialised, Latinate, or French-influenced forms. Such
items are difficult for rule-based systems and neural systems alike because the
correct pronunciation depends on etymological and lexical information not
recoverable from a small character-level training set.

This also clarifies why architecture changes have limited leverage. GQA, RoPE,
SwiGLU, and checkpoint averaging improve optimisation and representation, but
they do not add the missing linguistic variables. The most promising future
improvements are therefore not simply larger beams or deeper layers, but
sources of information aligned with the observed errors: stress prediction,
morphological decomposition, phonologically structured smoothing, or pretraining
on a larger pronunciation lexicon before fine-tuning on the coursework data.

The multilingual results therefore refine rather than overturn the orthographic
depth argument. The shared model transfers useful low-level evidence about
sequence transduction and recurrent grapheme--phone mappings, especially for
Spanish. However, language tags and shared parameters do not provide the missing
English stress and lexical information, so English remains the lower-resource
generalisation bottleneck.

\section{Related Work}
\label{sec:related}

Neural G2P has long been framed as sequence transduction. \citet{yao2015sequence}
showed that sequence-to-sequence LSTMs could rival strong conventional systems.
Transformer G2P systems were later shown to be effective by
\citet{yolchuyeva2020transformer}. The SIGMORPHON 2020 shared task formalised
multilingual G2P evaluation and showed that strong neural baselines remain
competitive, while releasing complete system outputs for error analysis
\citep{gorman2020sigmorphon}. Recent surveys place G2P within broader speech
technology pipelines and contrast finite-state, neural, and pretrained
approaches \citep{cheng2024survey}. Pretrained systems such as T5- or
ByT5-based G2P are attractive, but they require substantially more data and
compute than assumed in this coursework setting.

\section{Conclusion}
\label{sec:conclusion}

The encoder--decoder Transformer is effective for Spanish and benefits from
multilingual training, but strict unseen English G2P remains difficult. The
main limitation is not merely architecture but data scale: English requires
stress, lexical, and etymological information that 2,400 examples do not
robustly identify. Future work should test pretraining on CMUDict before
SIGMORPHON fine-tuning, phonologically informed label smoothing, and
morphology-aware augmentation for derived and compound words.

\bibliography{references}

\end{document}
